\subsection{Finite wall tests and the finite-rank interface}
\label{sec:flag-finite-wall}
The exact finite-cutoff comparison does not require the quantitative
wall approximation \texttt{ThmW} or coefficient error
\(O(\log B/B)\) of \texttt{ThmBA}. For this purpose it suffices
to establish qualitative convergence at finitely many fixed
rational weights on each cross-block wall. The number and choice
of these tests may depend on the fixed tower. They are not
new samples of the moving prime-sum parameter.

\paragraph{Exact comparison from qualitative wall limits.}
Fix the actual tower and write \(J_B=I_B/v\) for its normalized
arithmetic intersection functional on degree \(d+1\).
Assume the exact \texttt{Prop6} relation on this tower, the
source Frobenius shape condition, within-block antisymmetry
\texttt{I8} for \(J_B-H\), the wall vanishing of \(H\) in
\texttt{I9}, and the transitive label symmetry \texttt{PropA}
of \(J_B\) and invariance of \(H\) under that group.
For each label \(\tau\), consider the cross wall
\(u_{\tau,p-1}=u_{\tau,p}\). Its relatively open dominant
part, with all other adjacent gaps strict, contains a
relatively open set of rational test weights of norm
less than \(1/2\). Require \(J_B(\ell_u^{d+1})\to0\) for
a fixed finite set of such tests whose evaluations determine
the restrictions of polynomials of degree at most \(d+1\)
to this wall. Such a set exists: evaluations on the rational
points of the open part separate polynomials on the wall
(by density and the polynomial identity principle); finite
dimensionality gives a finite spanning subset of evaluations.

Here is a proof of the reduction. By (FA1)--(FA3), \(R_B\)
from (FA7) is a well-defined linear functional, and \(R_B\to0\)
coefficientwise on the fixed finite-dimensional space, by
absolute convergence for \(k\ge2\). Let
\[
 S=\sum_\tau\mathcal S_\tau,\qquad T=\sum_\tau\mathcal T_\tau
\]
be the two shift functionals of the tower-intersection reduction.
In the exact \texttt{Prop6} change, the non-tail contribution
is in \(\operatorname{span}(S,T)\). Indeed the subtracted
rank term in \texttt{deltaTerm} is
\(-(\log\ell)/(\ell-1)\) times the total first-block derivative.
The degree-one geometric-series part has coefficient
\[
 \kappa_\ell=\frac{\log\ell}{N}\sum_{a\ge1}\ell^{-a}
                 \mathrm{lieRank}(\mathrm{frob}(\ell,\tau,a))
\]
times the total derivative, independent of \(\tau\) by
the Frobenius shape. For \(k\ge2\) the total derivative is
zero on \(p_{\tau k}Q_{\tau k}\), since multiplication by
a positive full power sum has zero alternating degree.
These are exactly the signs in
\texttt{PowerSumDeltaDecomposition}. Extending by (FA1), the
change from \(B_0\) to \(B\) is
\[
 J_{B_0}-J_B
  =\sum_{B_0<\ell\le B,\ {\rm good\ prime}}
       \left(-\frac{\log\ell}{\ell-1}S+\kappa_\ell T\right)
       +R_{B_0}-R_B .
\]
Therefore the class of \(J_B-H-R_B\) modulo
\(\operatorname{span}(S,T)\) is independent of \(B\).

To see that it is zero, identify a functional \(L\) with the
homogeneous polynomial \(L(\ell_u^{d+1})\); this is an isomorphism
by the multinomial formula. The polynomials of \(J_B-H\)
are alternating inside the two blocks, in each label, by
\texttt{I8}. They tend to zero at the chosen cross-wall tests.
Inside the fixed finite-dimensional space of within-block
alternating polynomials, the common kernel of these evaluations
is the full Vandermonde times a block-invariant linear polynomial.
In fact the evaluations force vanishing on the entire indicated
wall; permutations within the blocks force vanishing on every
cross wall. The within-block Vandermonde factors were already
present. Divisibility by all distinct linear differences then
gives the assertion, since the total degree is \(d+1\).
The remaining linear polynomial is constant-coefficient on
each block. Intersect with the transitive label-invariant
subspace supplied by \texttt{PropA} and invariance of \(H\).
The coefficients are then common across labels.
This two-dimensional kernel is exactly \(\operatorname{span}(S,T)\)
under the isomorphism, by the exact shift formulas on
\(\ell_u^{d+1}\) and the flag degree product formula.
On a fixed complement of the kernel the evaluation map has
a linear left inverse, so the complement coefficients of
\(J_B-H\) tend to zero. Since \(R_B\to0\), the constant quotient
class identified above is zero. For each fixed cutoff \(B\)
there consequently exist \(A_B,b_B\) with
\[
 J_B=H-A_B S-b_B T+R_B .
\]
This is the exact source \texttt{IdentAt} and its extension
to all degree-\(d+1\) polynomials. It avoids estimating the
interpolation inverse as the field or dimension grows.
The assertion is conditional on the actual within-block
antisymmetry, label symmetry, wall convergence and two-cutoff
linkage; this argument does not establish those geometric
premises. In particular the determinant test alone does not
prove wall convergence.

\paragraph{Finite certificate for the lower image degree.}
For a projective normal variety \(X\) of dimension \(D\)
over the coefficient number field, let \(A\) be the original nef
weight line, with actual birational cusp-section system
\(V\subset H^0(X,A)\); include the cusp boundary conditions
in this system. On the normalization of its base-ideal blow-up
\(\pi:Y\to X\), write \(L=\pi^*A-E\), where \(E\) is the
effective Cartier divisor of the actual base ideal.
Then \(L\) is basepoint free and \(L^D=\delta\), the degree
of the actual cusp-section image. Let \(A'\) be the modified
nef weight line (the source \(u'\)).
It is sufficient to find a section \(\theta\) of
\(a(A-A')\), for some positive integer \(a\), whose pullback
satisfies \(\operatorname{div}(\pi^*\theta)\ge aE\).
Only the finitely many prime divisors in \(\operatorname{Supp}(E)\)
need to be tested; elsewhere the divisor of the section is
already effective. The valuations, including exceptional ones,
must be those of the normalized blow-up of the actual system,
not just the boundary primes of \(X\).
The resulting section exhibits \(a(L-\pi^*A')\) as effective.
Since \(L\) and \(\pi^*A'\) are nef, intersect this effective
divisor with
\(\sum_{i=0}^{D-1}L^i(\pi^*A')^{D-1-i}\) to deduce
\[
 \delta=L^D\ge(A')^D .
 \tag{FW1}
\]
The argument uses intersection positivity, not an asymptotic
rank threshold. No all-valuations test is necessary for this
sufficient certificate for the degree comparison.

For the actual weight shift \(u'_{\tau i}=u_{\tau i}
+3N{\bf1}_{i\ge p}\), provided \(u'\) remains strictly
dominant, the exact proportionality and Vandermonde formula
give the degree ratio
\[
 \frac{(A')^D}{A^D}
  =\prod_{\tau,a<p,b\ge p}\left(1-\frac{3N}{u_{\tau a}-u_{\tau b}}\right).
 \tag{FW2}
\]
All within-block differences cancel. Thus, if the divisor
certificate and actual top-degree proportionality hold,
the image loss \(\rho=1-\delta/A^D\) is bounded by
\[
 0\le\rho\le1-\prod_{\tau,a<p,b\ge p}
 \left(1-\frac{3N}{u_{\tau a}-u_{\tau b}}\right)
 \le 3N\sum_{\tau,a<p,b\ge p}\frac1{u_{\tau a}-u_{\tau b}}.
\]
The nonnegative lower bound uses \(L\le\pi^*A\) with both nef.
This retains the entire cross-gap distribution, rather than
only its minimum. In the finite-jet supplier already proved,
take any \(0\le\rho_*<1\) majorizing this loss. The actual
jet-count inequality gives, in its notation,
\[
 \frac{d_{\rm sec}}{d_{\rm lower}}\le
 \frac{\binom{\lfloor(t+D+3)/(1-\rho_*)\rfloor+D}{D}}
      {\binom{t-1}{D}}
\]
when \(d_{\rm lower}=\delta\binom{t-1}{D}\) and the other
stated weight/section hypotheses of that supplier hold.
Here the normalizations and section lattice have not been
changed. This is a finite sufficient rank estimate; it does
not construct \(\theta\), whose exceptional-divisor
vanishing is still a substantive obligation. Merely knowing
that the base locus is supported on the cusp boundary
does not give the required multiplicity.
